\section{Lo que se midió después}\label{sec:posteriores}

La evaluación final deja preguntas que sus datos no pueden responder. Dos se midieron aparte, con el
detector ya congelado y sin ajustar nada con lo que salió: qué pasa si el atacante reacciona a las
rotaciones, y qué pasa con pedazos más cortos que los que el contador cuenta.

\subsection{El atacante que se adapta}\label{sec:adaptativo}

\paragraph{La pregunta.} Una grabación, después de una rotación, sigue enviando la clave vieja. Un
atacante de verdad volvería a leer el archivo y seguiría con la nueva. La regla 5 del protocolo pedía
medir eso por separado.

\paragraph{El método.} \cod{python -m eval.final.adaptive} toma cada ataque grabado y lo reescribe
mientras lo reproduce:
\begin{itemize}
  \item lleva el valor vigente de cada clave, que cambia cada vez que el detector rota;
  \item antes de cada envío inserta una relectura del archivo de claves, con el valor vigente;
  \item cada envío manda las mismas posiciones, por el mismo canal y con la misma codificación, pero
    del valor \emph{vigente}.
\end{itemize}
La medición es exacta: por clave y por versión, las posiciones cuyos fragmentos salieron en llamadas
que el detector permitió. El ataque triunfa si de alguna versión saca más de 96 bits. Los ataques que
no se pueden reescribir de forma exacta se excluyen y se cuentan aparte; son los 10 de validación que
usan una tabla de sustitución propia.

\begin{table}[ht]
\caption{Atacante que se adapta, simulado (\cod{results/adaptive/summary.json}). Usa las marcas de
las grabaciones. «Escalado» es el porcentaje de ataques en que el detector llegó al tope de
rotaciones y pasó a bloquear.}\label{tab:adaptativo}
\footnotesize
\setlength{\tabcolsep}{4.5pt}
\begin{tabular}{@{}llrrlrrrr@{}}
\toprule
 & & & No & & \multicolumn{2}{c}{Bits} & Rotaciones & Escalado \\
\cmidrule(lr){6-7}
Set & Sistema & Reescritos & reescr. & Frenados & mediana & máximo & (mediana) & (\unit{\percent}) \\
\midrule
\tabladatos{tab_adaptativo}
\bottomrule
\end{tabular}
\end{table}

\begin{table}[ht]
\caption{El mismo análisis sobre el set de prueba, por familia: ataques frenados de 30 y máximo de
bits de una versión.}\label{tab:adaptativo-familia}
\small
\begin{tabular}{@{}lrrrrrr@{}}
\toprule
 & \multicolumn{2}{c}{Alicanto} & \multicolumn{2}{c}{Sin evidencia} & \multicolumn{2}{c}{Sin rotación} \\
\cmidrule(lr){2-3}\cmidrule(lr){4-5}\cmidrule(l){6-7}
Familia & frenados & máx. bits & frenados & máx. bits & frenados & máx. bits \\
\midrule
\tabladatos{tab_adaptativo_familia}
\bottomrule
\end{tabular}
\end{table}

\paragraph{Qué dice.} Con las marcas de las grabaciones, 180 de 180 ataques que se adaptan no pasan
de 96 bits de ninguna versión de la clave: el máximo observado es 64 y la mediana, 20.

\paragraph{Cómo lo logra, y a qué costo.} No lo logra rotando sin fin. En todos los ataques del set
de prueba el detector llega a su tope de tres rotaciones y entra en modo escalado: bloquea toda salida
de esa clave hacia destinos que no son los suyos. En total bloquea 831 envíos después de escalar. Es
una detención, no una rotación más, y tiene un costo: mientras dura, el agente no puede enviar nada
hacia afuera que no sea a los destinos de la clave. Ese costo sobre trabajo normal no está medido.

\begin{limite}
\textbf{Tres salvedades.}
\begin{itemize}
  \item Es una simulación. El atacante es un guion que relee y reenvía; no busca caminos nuevos.
  \item Usa las marcas de las grabaciones, igual que la corrida única. Con las marcas del proxy las
    familias A9 y A13 se comportarían distinto; ese cruce no se midió.
  \item En la ablación sin evidencia, la familia A9 queda exactamente en 96 bits y se cuenta como
    frenada. El servidor del atacante de la demo, con 96 bits, daría la clave por completada
    (\cref{sec:planteamiento}). No toca a las cifras de Alicanto.
\end{itemize}
\end{limite}

\subsection{Pedazos más cortos que el mínimo}\label{sec:fragmentos}

\paragraph{La pregunta.} La invariante de la \cref{sec:invariante} se demuestra para pedazos de 7
caracteres. El contador no cuenta los de 1 a 3. ¿Qué pasa con ellos?

\paragraph{El método.} \cod{python -m eval.final.short_fragments} construye sesiones sintéticas de
600 llamadas: una lectura del archivo de claves y, después, un envío cada cierto número de llamadas,
con pedazos que recorren la clave de corrido. Quien envía conoce las reglas y relee la clave en cuanto
rota, de modo que cada versión se mide en el peor caso. Se cruzan cinco tamaños de pedazo (1, 2, 3, 4
y 7 caracteres), tres ritmos (un envío en cada llamada, cada 8 y cada 51) y dos destinos, contra tres
configuraciones: 90 combinaciones. Las marcas son las del proxy.

\begin{table}[ht]
\caption{Máximo de bits de una misma versión de la clave que salió en cada combinación
(\cod{results/short_fragments/summary.json}). En negrita, las que pasan de 96. El destino no cambia
ningún valor. El \apref{ap:fragmentos} trae la tabla completa.}\label{tab:fragmentos}
\small
\begin{tabular}{@{}lrrrrrr@{}}
\toprule
 & Un envío & \multicolumn{5}{c}{Caracteres por envío} \\
\cmidrule(l){3-7}
Configuración & cada & 1 & 2 & 3 & 4 & 7 \\
\midrule
\tabladatos{tab_fragmentos_matriz}
\bottomrule
\end{tabular}
\end{table}

\paragraph{Qué dice.}
\begin{itemize}
  \item \textbf{Pedazos de 4 y de 7, los que el contador cuenta.} Ninguna de las 36 combinaciones
    pasa de 96 bits. El máximo, 84 bits, es el caso exacto de la invariante: 7 caracteres cada 51
    llamadas, con la configuración de la demo.
  \item \textbf{Pedazos de 1 a 3, con la configuración evaluada.} A lo más 12 bits. El contador no ve
    nada, pero la evidencia rota en el primer envío y, tras tres rotaciones, el modo escalado bloquea:
    al ritmo de un envío por llamada, de 590 envíos quedan bloqueados 587.
  \item \textbf{Pedazos de 1 a 3, sin la evidencia.} Aquí está el hueco. Con los parámetros
    congelados y sin modelo pasan de 96 bits 8 de 18 combinaciones; con la configuración de la demo y
    del \esp{hook}, 10 de 18. Un carácter por llamada saca la clave entera antes de la primera
    rotación periódica.
\end{itemize}

\begin{clave}
\textbf{La consecuencia práctica.} Contra pedazos muy cortos lo que protege es la evidencia de
comportamiento, no el contador ni la invariante. La configuración de la demo y la del \esp{hook} de
Claude Code no la llevan. Por eso el informe dice que el \esp{hook}, tal como está, no cubre este
caso.
\end{clave}

\subsection{Las ablaciones sobre validación}

Después de la evaluación, una investigación midió sobre validación, sin tocar el detector y sin volver
a usar el set de prueba, qué aporta cada componente. Sus conclusiones están en
\cod{results/test/LECTURA.md} y las usamos en la \cref{sec:evidencia} y en la \cref{sec:no-conectado}:
las señales de los grafos reducen las rotaciones en sesiones normales pero no cambian los ataques
frenados; el modelo equivale a una regla simple; y la propagación de sospecha entre pasos lejanos no
está conectada a las decisiones.

\subsection{Los modelos reales}

Las corridas con modelos reales están en la \cref{sec:llm}. Se resumen en una línea: en 21 corridas
por brazo ningún modelo siguió la instrucción plantada, así que no hubo fuga que frenar.
